\chapter{Caractéristiques fondamentales, prévisions des analystes et événements}\label{ch:rs:fundamental-analyst-and-event-features}

Une société clôture son trimestre le 31 décembre, annonce ses résultats le 25 janvier et dépose son rapport trimestriel le 5
février. Une base de données indexée sur la fin de l'exercice affiche les résultats du trimestre au 31 décembre, et un backtest construit sur elle opère
sur l'annonce trois semaines et demie avant que quiconque ait pu la lire. L'erreur est invisible dans le code, qui se joint
sur une colonne de dates comme sur n'importe quelle autre, et spectaculaire dans les résultats. Ce chapitre indexe les chiffres comptables, les surprises, les prévisions des
analystes et les dates d'événements sur l'instant où ils sont devenus connus. Ses données réelles sont les résultats publiés de dix grandes sociétés américaines et
les calendriers de dépôt tirés des interfaces publiques EDGAR de la SEC~; sa mesure des dégâts porte sur \texttt{firm.synthmkt}, où la
vérité est connue.

\section{Données comptables, point-in-time}

\begin{definition}[Exercice comptable, délai de publication]\label{def:rs:fundamental-analyst-and-event-features:fiscal}
Un \emph{exercice comptable}\index{exercice comptable} est le trimestre ou l'année que décrit un rapport comptable, identifié par ses dates de début
et de fin~; le \omterm{def:rs:point-in-time-data-and-the-biases:bitemporal}{temps de validité} de chaque chiffre du rapport est cet exercice. Le \emph{délai de publication}\index{délai de publication} est le temps écoulé entre
la fin de l'exercice et le dépôt du rapport auprès du régulateur.
\end{definition}

Un chiffre comptable possède au moins trois dates après la fin de son exercice~: la publication des résultats (un communiqué de presse, déposé aux États-Unis sous la
forme d'un Form 8-K au titre de l'item 2.02), le rapport périodique (Form 10-Q pour un trimestre, 10-K pour une année), et tout dépôt ultérieur qui reprend le
même exercice à titre de comparatif, éventuellement avec une valeur différente
(la \cref{fig:rs:fundamental-analyst-and-event-features:timeline}). Le marché apprend le chiffre à la publication~; une base construite à partir des rapports périodiques
l'apprend au dépôt~; une base qui ne stocke que la fin de l'exercice prétend le connaître dès cette fin.

\begin{omfigure}
\begin{tikzpicture}[font=\small,x=0.19cm]
\draw[thick,-Stealth] (0,0) -- (66,0) node[right]{jours};
\foreach \x/\l in {0/0,23/23,29/29,40/40,45/45,60/60} {\draw (\x,0.12) -- (\x,-0.12) node[below]{\l};}
\node[above,align=center,text=omThm] at (0,0.35) {fin de\\exercice};
\node[above,align=center,text=omDef] at (23,0.35) {publication des\\résultats};
\node[above,align=center,text=omDef] at (31,1.35) {10-Q\\déposé};
\draw[omDef] (29,0.15) -- (30,1.3);
\draw[black!50,dashed] (40,-0.6) -- (40,2.3) node[above,align=center,font=\footnotesize,text=black!70]{échéance du 10-Q,\\grands déposants};
\draw[black!50,dashed] (45,-0.6) -- (45,1.2) node[above,align=center,font=\footnotesize,text=black!70]{autres};
\draw[black!50,dashed] (60,-0.6) -- (60,1.2) node[above,align=center,font=\footnotesize,text=black!70]{10-K,\\grands déposants};
\draw[omThm,thick,Stealth-Stealth] (0,-1.1) -- node[below,font=\footnotesize]{avance d'une clé indexée sur la fin de l'exercice} (23,-1.1);
\end{tikzpicture}

\omcaption{Les dates d'un chiffre trimestriel, aux médianes de dix grandes sociétés américaines pour leurs trois premiers exercices trimestriels depuis 2013 (23 jours avant
la publication des résultats, 29 jours avant le 10-Q), avec les échéances de dépôt de la SEC.}\label{fig:rs:fundamental-analyst-and-event-features:timeline}
\end{omfigure}

\begin{dated}{2026-09}{Échéances de dépôt périodique aux États-Unis}\label{dat:rs:fundamental-analyst-and-event-features:deadlines}
Les instructions de la SEC relatives au Form 10-Q exigent le rapport trimestriel dans les 40 jours suivant la fin du trimestre pour les grands déposants accélérés et les déposants accélérés
et dans les 45 jours pour tous les autres déclarants~; aucun 10-Q n'est déposé pour le quatrième trimestre. Le Form 10-K est dû 60 jours après la fin de l'exercice
pour les grands déposants accélérés, 75 pour les déposants accélérés et 90 pour tous les autres. Les interfaces de programmation d'applications EDGAR de la SEC mettent à disposition
l'historique de dépôts de chaque déclarant et les données XBRL de ses états financiers au format JSON, mises à jour en temps réel, sans authentification ni clés.
\end{dated}

Les dix sociétés de \texttt{rs\_fetch\_edgar.py} (Apple, Microsoft, Walmart, Coca-Cola, Johnson \& Johnson, Exxon Mobil,
Procter \& Gamble, Intel, Home Depot et Nike) comptent 542 fins d'exercice depuis 2013 dans leurs historiques de dépôts EDGAR, dont 541
suivies d'une publication des résultats dans les 90 jours. Pour les trois premiers trimestres, la société médiane publiait ses résultats 23 jours
après la fin de l'exercice et déposait son 10-Q 29 jours après~; pour les exercices annuels, les médianes sont de 26 et 48 jours. Le plus long délai du 10-Q est
de 41 jours et le plus long délai du 10-K de 60, dans les échéances. Seuls 16~\% des rapports ont été déposés le jour de la publication~; l'écart médian est
une semaine, pendant laquelle le chiffre est public mais absent d'une base alimentée par les dépôts
(la \cref{fig:rs:fundamental-analyst-and-event-features:calendar}).

\begin{omfigure}
\begin{tikzpicture}
\begin{axis}[width=11cm,height=6.2cm,xbar,bar width=5pt,xlabel={médiane des jours après la fin du trimestre},
  tick label style={font=\small},label style={font=\small},xmin=0,xmax=42,ymin=-0.6,ymax=9.6,ytick={0,1,2,3,4,5,6,7,8,9},
  yticklabels={Home Depot,Johnson \& Johnson,Walmart,Nike,Procter \& Gamble,Coca-Cola,Intel,Microsoft,Exxon Mobil,Apple},
  y dir=reverse,grid=major,grid style={gray!25},legend columns=2,
  legend style={font=\footnotesize,at={(0.5,-0.2)},anchor=north,draw=none,fill=none,
  /tikz/every even column/.append style={column sep=8pt}},table/col sep=comma]
\addplot[fill=omDef!55,draw=omDef] table[x=release,y=k]{figdata/research/11-fundamental-analyst-and-event-features/calendar.csv};
\addplot[fill=gray!35,draw=gray] table[x=filing,y=k]{figdata/research/11-fundamental-analyst-and-event-features/calendar.csv};
\draw[black!60,dashed] (axis cs:40,-0.6) -- (axis cs:40,9.6);
\legend{publication des résultats (8-K 2.02),10-Q déposé}
\end{axis}
\end{tikzpicture}

\omcaption{Médiane des jours écoulés entre la fin de chacun des trois premiers exercices trimestriels et la publication des résultats, puis le dépôt du 10-Q, dix
sociétés, 2013 à 2026~; la ligne en pointillés marque l'échéance de 40 jours. Certaines déposent avec la publication, d'autres quelques semaines plus tard. Données~: SEC EDGAR
submissions, via \texttt{rs\_fetch\_edgar.py}.}\label{fig:rs:fundamental-analyst-and-event-features:calendar}
\end{omfigure}

Les chiffres changent aussi après leur premier dépôt. Chaque fait XBRL dans EDGAR liste chaque dépôt qui l'a rapporté, de sorte que le
bénéfice dilué par action d'un trimestre peut être suivi depuis son premier rapport jusqu'à chaque comparatif ultérieur. Parmi les dix sociétés, 31
valeurs trimestrielles ont changé~:

\begin{center}
\small
\begin{tabular}{@{}lccl@{}}
\toprule
entreprise & valeurs modifiées & par un ratio exact & ratios\\
\midrule
Apple & 13 & 13 & 7 et 4\\
Nike & 7 & 6 & 2\\
Microsoft & 5 & 0 & ---\\
Walmart & 3 & 3 & 3\\
Coca-Cola & 2 & 2 & 2\\
Johnson \& Johnson & 1 & 0 & ---\\
les quatre autres & 0 & 0 & ---\\
\bottomrule
\end{tabular}
\end{center}

Vingt-quatre des 31 valeurs modifiées l'ont été selon un ratio exact (au centime près), signature d'une division d'actions~: les 10-K d'Apple mentionnent une division par sept
le 6 juin 2014 et une division par quatre le 28 août 2020, après quoi toutes les informations par action ont été rétroactivement ajustées. Les autres
sept sont des \omterm{def:rs:point-in-time-data-and-the-biases:vintage}{retraitements} de fond. Le bénéfice dilué de Microsoft pour son exercice 2017 a d'abord été publié à 2.71~\$ et un an
plus tard à 3.25~\$~: elle a adopté une nouvelle norme de chiffre d'affaires dans son exercice 2018 selon la méthode rétrospective complète, qui l'obligeait à retraiter
chacun des exercices antérieurs présentés. Un backtest de l'exercice 2017 doit utiliser 2.71~\$ jusqu'en août 2018.

\begin{definition}[Douze derniers mois]\label{def:rs:fundamental-analyst-and-event-features:ttm}
Le \emph{douze derniers mois}\index{douze derniers mois} (TTM) d'un flux (résultats, chiffre d'affaires, flux de trésorerie) à la fin d'un trimestre est
la somme des quatre trimestres consécutifs les plus récents, chacun tel qu'il est connu à l'instant de décision~; un quatrième trimestre qui n'est publié seulement à l'intérieur
d'un chiffre annuel est l'année moins les neuf premiers mois.
\end{definition}

Additionner quatre trimestres est l'endroit où les dates font mal. Les bénéfices dilués trimestriels publiés en premier par Apple pour son exercice 2014 étaient de 14.50~\$ et
11.62~\$ (avant la division de juin) et 1.28~\$ et 1.42~\$ (après)~: un TTM construit à partir des premières publications vaut 28.82~\$, alors que le chiffre annuel d'Apple
publié pour le même exercice était de 6.45~\$. Le même phénomène se produit à travers la division de 2020, avec 10.85~\$ contre un chiffre annuel de 3.28~\$
(la \cref{fig:rs:fundamental-analyst-and-event-features:ttm}). La vue d'aujourd'hui pose le problème inverse~: elle réajuste l'historique en connaissant
les divisions postérieures, et elle est elle-même un patchwork, car une période n'est réajustée que si un dépôt ultérieur l'a de nouveau publiée. La règle consiste à
conserver les valeurs par action telles que publiées en premier, avec leurs \omterm{def:rs:point-in-time-data-and-the-biases:bitemporal}{temps de connaissance}, conserver les divisions comme des événements avec les leurs (chapitre~4), et réajuster au moment
de la requête vers la base par action à la date de décision, soit la base du cours par lequel le nombre sera divisé.

\begin{omfigure}
\begin{tikzpicture}
\begin{semilogyaxis}[width=11cm,height=5.6cm,xlabel={fin de trimestre},ylabel={bénéfice dilué TTM (\$)},tick label style={font=\small},
  label style={font=\small},xmin=2011.5,xmax=2026.7,ymin=2,ymax=80,ytick={2,5,10,20,50},yticklabels={2,5,10,20,50},
  grid=major,grid style={gray!25},
  xticklabel style={/pgf/number format/1000 sep=},legend pos=north east,legend style={font=\footnotesize},table/col sep=comma]
\addplot[omThm,thick,mark=*,mark size=1.2pt] table[x=year,y=first]{figdata/research/11-fundamental-analyst-and-event-features/apple_ttm.csv};
\addplot[omDef,thick,dashed] table[x=year,y=latest]{figdata/research/11-fundamental-analyst-and-event-features/apple_ttm.csv};
\draw[black!50] (axis cs:2014.43,2) -- (axis cs:2014.43,80) node[pos=0.06,right,font=\footnotesize]{7:1};
\draw[black!50] (axis cs:2020.66,2) -- (axis cs:2020.66,80) node[pos=0.06,right,font=\footnotesize]{4:1};
\legend{à partir des trimestres publiés en premier,à partir des valeurs d'aujourd'hui}
\end{semilogyaxis}
\end{tikzpicture}

\omcaption{Bénéfice dilué d'Apple sur les \omterm{def:rs:fundamental-analyst-and-event-features:ttm}{douze derniers mois} à la fin de chaque trimestre, additionné à partir de la valeur publiée en premier pour chaque trimestre et à partir de
sa dernière valeur dans EDGAR, avec les deux divisions d'actions. Aucune des deux n'est une série exploitable sans ajustement des divisions à la date de décision. Données~: SEC
EDGAR company facts.}\label{fig:rs:fundamental-analyst-and-event-features:ttm}
\end{omfigure}

\section{Ratios comme prédicteurs}

\begin{definition}[Ratio comptable]\label{def:rs:fundamental-analyst-and-event-features:ratio}
Un \emph{ratio comptable}\index{ratio comptable} divise une grandeur comptable par une autre ou par une grandeur de marché~: book-to-price,
rendement des résultats (résultats TTM sur cours), chiffre d'affaires sur cours, rentabilité des capitaux propres, effet de levier, croissance des actifs.
\end{definition}

Un ratio a deux \omterm{def:rs:point-in-time-data-and-the-biases:bitemporal}{temps de connaissance} et doit respecter les deux. Le côté comptable est connu à son dépôt (ou à sa publication)~; le côté marché
(cours, capitalisation boursière) est connu à son propre instant, normalement l'instant de décision. Les appariements qui échouent sont réguliers~: un bénéfice TTM
rapporté à un cours sur une autre base par action~; une valeur comptable indexée sur sa fin de période~; un ratio dont le dénominateur peut changer de signe
(résultats, capitaux propres comptables), ce qui exige un traitement séparé pour les valeurs négatives, plutôt qu'un classement qui place le rendement des résultats le plus
négatif à côté de l'action la plus chère. Le book-to-price dans \texttt{firm.synthmkt} montre à quel point la date de référence compte pour une variable lente~:
son IC moyen de rang sur les 21 jours suivants vaut 0.029 (erreur type 0.038) indexé sur la fin de période et $-0.005$ (0.023) indexé sur le jour du
dépôt. La différence est du bruit~: chaque IC trimestriel est dominé par le rendement réalisé du facteur valeur sur sa fenêtre, et déplacer la fenêtre
de quelques semaines modifie peu l'espérance, parce que rien sur le marché simulé ne se produit le jour du dépôt. Le dommage causé par une mauvaise
date de référence est proportionnel à l'ampleur avec laquelle la publication du chiffre fait bouger le cours, et la publication d'une valeur comptable ne le fait guère bouger.

\section{Surprises et révisions}

\begin{definition}[Surprise de résultats, résultat inattendu standardisé]\label{def:rs:fundamental-analyst-and-event-features:sue}
Une \emph{surprise de résultats}\index{surprise de résultats} est le résultat publié d'une période moins sa valeur espérée~: une prévision consensuelle
des analystes (livre~2, chapitre~31) ou un modèle de séries temporelles. Le \emph{résultat inattendu standardisé}\index{résultat inattendu standardisé} (SUE)
divise la surprise d'une marche aléatoire saisonnière, $x_q - x_{q-4}$, par l'écart type de cette différence sur les huit
trimestres précédents.
\end{definition}

Foster, Olsen et Shevlin (1984) ont établi que les dérives systématiques des rendements après annonce sont associées au signe et à l'ampleur de
résultats inattendus, et que, pour les modèles d'espérance fondés sur la série temporelle des résultats trimestriels publiés, l'erreur de prévision et la taille
de l'entreprise expliquaient 81~\% et 61~\% de la variation des dérives. La publication est l'événement~; la dérive qui suit est l'opportunité du prédicteur~; une
date de référence antérieure à la publication capture la publication elle-même. Dans \texttt{firm.synthmkt}, les résultats de chaque trimestre sont annoncés un nombre aléatoire de jours de bourse après sa fin (31 en moyenne, soit environ six semaines~: l'avance d'une date de référence en fin de période), avec un saut dans le sens de la vraie
surprise et une dérive plus faible ensuite. La surprise indexée sur son annonce a un IC moyen de rang avec le rendement corrigé du
marché sur les 21 jours suivants de 0.020 (erreur type 0.006)~: la dérive introduite. Indexée sur la fin de période, elle vaut 0.167, huit fois plus élevée, parce qu'un tiers des
fenêtres contiennent le saut de l'annonce (la \cref{fig:rs:fundamental-analyst-and-event-features:ic}).

\begin{omfigure}
\begin{tikzpicture}
\begin{axis}[width=10cm,height=5.4cm,ybar,bar width=16pt,ylabel={IC de rang moyen, 21 jours suivants},tick label style={font=\small},
  label style={font=\small},xtick={0,1,2,3},xticklabels={{SUE\\connu},{SUE\\fin de période},{B/P\\connu},{B/P\\fin de période}},
  xticklabel style={align=center},
  ymin=-0.1,ymax=0.2,enlarge x limits=0.18,grid=major,grid style={gray!25},table/col sep=comma]
\addplot[fill=omDef!55,draw=omDef,error bars/.cd,y dir=both,y explicit]
  table[x=k,y=ic,y error=se2]{figdata/research/11-fundamental-analyst-and-event-features/ic.csv};
\end{axis}
\end{tikzpicture}

\omcaption{IC de rang moyen sur 39 coupes transversales trimestrielles de \texttt{firm.synthmkt} (barres~: deux erreurs types) de la vraie \omterm{def:rs:fundamental-analyst-and-event-features:sue}{surprise de résultats}
et du book-to-price, indexés sur le jour où chacun est devenu connu (annonce, dépôt) et sur la fin de la période. La date de référence en fin de période multiplie
par huit l'IC de la surprise et laisse le book-to-price dans son bruit.}\label{fig:rs:fundamental-analyst-and-event-features:ic}
\end{omfigure}

\begin{definition}[Révision des analystes, dispersion des prévisions]\label{def:rs:fundamental-analyst-and-event-features:revision}
Une \emph{révision des analystes}\index{révision des analystes} est une modification de la prévision ou de la recommandation d'un analyste sur un titre~; en tant que prédicteur, la
variation de la prévision consensuelle sur une fenêtre, rapportée au cours ou à la dispersion des prévisions. \emph{Dispersion
des prévisions}\index{dispersion des prévisions} est l'écart type entre analystes des prévisions courantes, divisé par la valeur absolue de leur
moyenne.
\end{definition}

Womack (1996) a établi que les nouvelles recommandations d'achat et de vente émises par des analystes de grandes maisons de courtage américaines font bouger les cours d'un coup, bien que peu coïncident
avec une nouvelle information publique, et que les réactions sont incomplètes~: la dérive moyenne après une recommandation d'achat est modeste (+2.4~\%) et brève,
après une recommandation de vente plus forte ($-9.1\%$) et d'une durée de six mois. Diether, Malloy et Scherbina (2002) ont établi que les actions à plus forte \omterm{def:rs:fundamental-analyst-and-event-features:revision}{dispersion
des prévisions} de résultats des analystes rapportent des rendements futurs plus faibles que des titres comparables, surtout parmi les petites valeurs et les perdants passés, ce qui s'accorde
avec des cours qui reflètent le point de vue des optimistes lorsque les investisseurs les plus pessimistes n'interviennent pas. Les deux sont d'abord des problèmes de \omterm{def:rs:point-in-time-data-and-the-biases:lookahead}{données point-in-time}~: une
base de prévisions doit conserver chaque prévision avec la date à laquelle l'analyste l'a faite et la date à laquelle le fournisseur l'a enregistrée, et une recommandation qui
est ensuite retirée ou corrigée doit rester dans l'historique tel qu'il était. La couverture constitue en soi un biais~: les analystes suivent les sociétés dont ils attendent de bonnes
performances et abandonnent les autres, de sorte que l'ensemble des sociétés couvertes évolue avec l'issue.

\section{Calendriers des événements}

\begin{definition}[Calendrier des événements, fenêtre d'annonce]\label{def:rs:fundamental-analyst-and-event-features:calendar}
Un \emph{calendrier des événements}\index{calendrier des événements} recense les événements d'entreprise programmés (publications de résultats, dates de dividende, assemblées d'actionnaires,
changements d'indices), avec la date à laquelle chacun est annoncé et la date où il survient. Une \emph{fenêtre d'annonce}\index{fenêtre d'annonce} est une
plage fixe de jours de bourse autour d'un événement, par exemple $[-1, +1]$, utilisée pour mesurer sa réaction ou pour l'exclure.
\end{definition}

Un calendrier est lui-même des \omterm{def:rs:point-in-time-data-and-the-biases:lookahead}{données point-in-time}~: une date programmée est annoncée avant l'événement, peut se déplacer et est conservée avec le \omterm{def:rs:point-in-time-data-and-the-biases:bitemporal}{temps
de connaissance} auquel elle a été annoncée~; une date qui se déplace vaut elle-même d'être testée comme prédicteur. Les \omterm{def:rs:fundamental-analyst-and-event-features:fiscal}{délais de publication} des dix sociétés sont réguliers~: Home Depot a publié
16 jours après la fin de chacun de ses trois premiers trimestres comptables depuis 2013, Exxon Mobil 30 jours à deux jours près, tandis que celui d'Apple varie davantage (un écart type de cinq jours). Une date prédite est donc un prédicteur exploitable avant que la
date confirmée n'existe, et \emph{jours avant le prochain événement} est une variable de conditionnement dont tout système de recherche a besoin~: elle
active des prédicteurs (une surprise après la publication), les désactive (une estimation de volatilité à cheval sur une publication), et dimensionne le risque autour de l'inconnu connu.
Les \omterm{def:rs:fundamental-analyst-and-event-features:calendar}{fenêtres d'annonce} servent à l'inverse~: un prédicteur de retournement ou de volatilité estimé à travers un jour de publication mesure l'annonce,
et non l'effet qu'il devait mesurer, si bien que le code de recherche exclut la fenêtre ou la modélise à part.

\section{Fiches de prédicteur}

\begin{predictorcard}{Résultat inattendu standardisé (SUE)}\label{pred:rs:fundamental-analyst-and-event-features:sue}
\sfield{Définition} $(x_q - x_{q-4})/\sigma_8$, la surprise saisonnière en marche aléatoire du BPA trimestriel rapportée à l'écart type des
huit derniers de ces écarts~; ou la surprise par rapport au consensus, rapportée à la dispersion.
\sfield{Entrées et horodatages} BPA trimestriel tel que publié à l'origine, indexé sur la publication (le 8-K) ou, à partir des seuls dépôts, sur le dépôt~; facteurs de division
connus à la date de décision.
\sfield{Justification} Dérive postannonce (Foster, Olsen et Shevlin, 1984).
\sfield{Horizon et demi-vie} De quelques semaines à un trimestre après la publication~; sur \texttt{firm.synthmkt}, IC~0.020 sur 21 jours avec la bonne clé.
\sfield{Normalisation} Rang parmi les actions ayant annoncé dans la fenêtre~; \omterm{def:rs:anatomy-of-a-predictor:neutralise}{neutralisation} de la taille.
\sfield{Modes de défaillance} Une clé de fin de période (IC~0.167 sur le marché simulé~: tout cela relève de l'annonce)~; une surprise mesurée contre un chiffre
retraité~; un consensus périmé.
\sfield{Sources} Telles que citées~; \texttt{firm.fundpit.sue}~; \texttt{rs\_fundamentals.ic\_inflation}.
\end{predictorcard}

\begin{predictorcard}{Book-to-price, tel que connu}\label{pred:rs:fundamental-analyst-and-event-features:bp}
\sfield{Définition} Les capitaux propres comptables les plus récents connus à l'instant de décision, rapportés à la capitalisation boursière à l'instant de décision.
\sfield{Entrées et horodatages} Capitaux propres comptables indexés sur leur dépôt~; nombre d'actions et cours à la date de décision, sur une même base d'actions.
\sfield{Justification} La prime de valeur (livre~8, chapitre~6).
\sfield{Horizon et demi-vie} De quelques mois à quelques années~; le bilan change chaque trimestre, le cours chaque jour.
\sfield{Normalisation} Rang au sein de l'industrie~; capitaux propres négatifs traités à part.
\sfield{Modes de défaillance} Une clé de fin de période (peu de dégâts pour un ratio lent~: sur le marché simulé, l'écart reste dans le bruit)~;
bases d'actions incohérentes~; valeurs comptables retraitées utilisées avant leur \omterm{def:rs:point-in-time-data-and-the-biases:vintage}{retraitement}.
\sfield{Sources} \texttt{rs\_fundamentals.ic\_inflation}.
\end{predictorcard}

\section{Tutoriel~: les cinq semaines d'avance}\label{tut:rs:fundamental-analyst-and-event-features:tut}

\begin{tutorial}
\textbf{Objectif.} Construire les dates des publications de résultats réelles depuis EDGAR, les stocker en point-in-time, dériver les trimestres et les TTM tels que connus à
une date quelconque, et mesurer sur un marché synthétique ce qu'une clé de fin de période fait à une surprise et à un ratio. \textbf{État final~:}
\cref{fig:rs:fundamental-analyst-and-event-features:calendar,fig:rs:fundamental-analyst-and-event-features:ttm,fig:rs:fundamental-analyst-and-event-features:ic}~;
le tableau des \omterm{def:rs:point-in-time-data-and-the-biases:vintage}{retraitements}.
\begin{tutsteps}
\item \textbf{Récupérer} avec \texttt{rs\_fetch\_edgar.py}~: les soumissions (8-K item 2.02, 10-Q, 10-K) et les faits de l'entreprise (BPA dilué de chaque dépôt
qui l'a publié), avec un User-Agent descriptif et une pause entre les requêtes.
\item \textbf{Stocker} chaque fait sous sa durée, valide à la fin de la période, connu au dépôt.
\omcode{code/firm/fundpit/firm_fundpit.py}{42}{54}{Faits versés dans un magasin point-in-time, par durée.}\label{lst:rs:fundamental-analyst-and-event-features:load}
\item \textbf{Trimestres tels que connus}, le quatrième dérivé de l'année et des neuf mois.
\omcode{code/firm/fundpit/firm_fundpit.py}{57}{72}{Valeurs trimestrielles telles que connues à une date.}\label{lst:rs:fundamental-analyst-and-event-features:quarterly}
\item \textbf{Exécuter} \texttt{calendar\_summary()}, \texttt{restatements()}, \texttt{apple\_ttm()}, \texttt{ic\_inflation()} et
\texttt{fig\_fundamentals.py}.
\end{tutsteps}
\textbf{Pour aller plus loin.} Indexer la surprise sur le dépôt du 10-Q plutôt que sur la publication et mesurer ce que coûte la semaine qui les sépare~;
ramener les premiers trimestres publiés par Apple à la base d'actions de chaque date de décision, puis reconstruire une série de TTM propre.
\end{tutorial}

\section{Construction~: le magasin de caractéristiques fondamentales}\label{bld:rs:fundamental-analyst-and-event-features:fundpit}

\begin{build}
\bfield{Rôle} Caractéristiques comptables pour la recherche sur l'entreprise miniature, qui ne puissent pas fuir~: chaque valeur porte sa période et son \omterm{def:rs:point-in-time-data-and-the-biases:bitemporal}{temps
de connaissance}, et chaque valeur dérivée (trimestres, TTM, SUE, ratios) est calculée à partir de ce qui était connu à la date de décision.
\bfield{Interface} \texttt{duration\_kind}, \texttt{load(facts, field)}, \texttt{quarterly(store, entity, field, known)}, \texttt{ttm},
\texttt{sue(values, n)}, \texttt{restated}, \texttt{split\_ratio}, \texttt{revision}, \texttt{dispersion}, \texttt{days\_to\_event},
\texttt{in\_window}~; le stockage est \texttt{firm.pit.Store}.
\bfield{Règles} Le \omterm{def:rs:point-in-time-data-and-the-biases:bitemporal}{temps de validité} est la fin de la période, le \omterm{def:rs:point-in-time-data-and-the-biases:bitemporal}{temps de connaissance} le dépôt (ou la publication lorsqu'elle est stockée)~; un quatrième trimestre n'existe que lorsque
l'année et les neuf mois sont tous deux connus~; une valeur par action n'est jamais comparée entre bases d'actions.
\bfield{Tests d'acceptation} \texttt{code/firm/fundpit/tests/}~: durées d'une année construite à la main, le quatrième trimestre (4.80~\$ moins 3.30~\$), un trimestre
inconnu avant son dépôt, TTM, la vue ultérieure après une division~; liste des \omterm{def:rs:point-in-time-data-and-the-biases:vintage}{retraitements} et ratios de division au centime près~; SUE d'un saut saisonnier et d'une
marche aléatoire~; révision, dispersion, jours avant le prochain événement et fenêtres vérifiés à la main.
\bfield{Pour aller plus loin} Dates de publication tirées des dépôts 8-K comme \omterm{def:rs:point-in-time-data-and-the-biases:bitemporal}{temps de connaissance} de chaque trimestre~; un magasin de consensus avec horodatages des analystes et des fournisseurs~;
les événements de division de \texttt{firm.secmaster} appliqués au moment de la requête.
\end{build}

\begin{omsources}
\item US Securities and Exchange Commission, Form 10-Q et Form 10-K, instructions générales~; interfaces de programmation applicative d'EDGAR
(data.sec.gov).
\item Apple Inc., Form 10-K pour l'exercice 2014 et l'exercice 2020~; Microsoft Corporation, Form 10-K pour l'exercice 2018 (SEC EDGAR).
\item G.~Foster, C.~Olsen et T.~Shevlin, ``Earnings releases, anomalies, and the behavior of security returns'', \emph{The Accounting Review}
59(4), 1984.
\item K.~L.~Womack, ``Do brokerage analysts' recommendations have investment value?'', \emph{Journal of Finance} 51(1), 1996.
\item K.~B.~Diether, C.~J.~Malloy et A.~Scherbina, ``Differences of opinion and the cross section of stock returns'', \emph{Journal of
Finance} 57(5), 2002.
\end{omsources}

\section{Exercices}

\begin{exercise}[$\star$]\label{exo:rs:fundamental-analyst-and-event-features:1}
L'\omterm{def:rs:fundamental-analyst-and-event-features:fiscal}{exercice comptable} d'une entreprise se termine le 30 juin. Elle publie un BPA dilué sur neuf mois de 3.30~\$ au 31 mars et un BPA annuel de 4.80~\$. Calculer le BPA
du quatrième trimestre, et dire à quelle date il est connu pour la première fois si le 10-K est déposé 55 jours après la fin de l'exercice~?
\end{exercise}

\begin{exercise}[$\star$]\label{exo:rs:fundamental-analyst-and-event-features:2}
Le trimestre d'un déposant accéléré se termine le 30 septembre. Déterminer à quelle date au plus tard son 10-Q doit être déposé, et la date limite du 10-K d'un grand déposant accéléré pour un exercice
se terminant le 31 décembre (année non bissextile).
\end{exercise}

\begin{exercise}[$\star$]\label{exo:rs:fundamental-analyst-and-event-features:3}
Le BPA d'un trimestre a d'abord été publié à 1.42~\$, puis à 0.355~\$. Expliquer ce qui s'est passé et comment le magasin doit conserver les deux valeurs.
\end{exercise}

\begin{exercise}[$\star\star$]\label{exo:rs:fundamental-analyst-and-event-features:4}
Les BPA trimestriels sur douze trimestres sont 1.00, 1.10, 0.90, 1.20, 1.04, 1.12, 0.96, 1.22, 1.06, 1.16, 0.94, 1.40. Calculer les écarts saisonniers et le SUE du dernier trimestre avec $n = 6$ (les six écarts précédents, écart type d'échantillon).
\end{exercise}

\begin{exercise}[$\star\star$]\label{exo:rs:fundamental-analyst-and-event-features:5}
Si les annonces tombent uniformément entre 0 et 62 jours de bourse après la fin d'un trimestre, déterminer la part des fenêtres de 21 jours commençant le lendemain de la
fin du trimestre qui contient l'annonce, et la mettre en relation avec l'inflation de l'IC d'un facteur huit.
\end{exercise}

\begin{exercise}[$\star\star$]\label{exo:rs:fundamental-analyst-and-event-features:6}
Expliquer pourquoi le book-to-price indexé sur la fin de période peut montrer peu d'inflation alors que la \omterm{def:rs:fundamental-analyst-and-event-features:sue}{surprise de résultats} en montre beaucoup,
et dire ce qui ferait aussi gonfler le book-to-price.
\end{exercise}

\begin{exercise}[$\star\star\star$]\label{exo:rs:fundamental-analyst-and-event-features:7}
\emph{Programmation.} À partir de \texttt{data/research/edgar\_filings.csv}, calculer pour chaque entreprise la part des trimestres dont le 10-Q a été déposé le
jour de la publication, et le nombre médian de jours entre la publication et le dépôt pour les autres.
\end{exercise}

\begin{exercise}[$\star\star\star$]\label{exo:rs:fundamental-analyst-and-event-features:8}
\emph{Trouver la faille.} «~Nous utilisons la table historique des fondamentaux du fournisseur, qui est mise à jour chaque fois qu'une entreprise procède à un \omterm{def:rs:point-in-time-data-and-the-biases:vintage}{retraitement}, si bien que notre historique de résultats
est toujours le plus exact disponible.~»
\end{exercise}

\section{Problème~: L'avance de cinq semaines}

\begin{problem}[{Problème du week-end --- une date clé, mesurée}]\label{pb:rs:fundamental-analyst-and-event-features:1}
Les données EDGAR de dix entreprises depuis 2013, et \texttt{firm.synthmkt} avec leurs annonces de résultats, dépôts et valeurs comptables.

\medskip\noindent\textbf{Partie I --- Les dates.}
\begin{enumerate}
\item Combien de fins de période dix entreprises ont-elles connues depuis 2013, et combien sont suivies d'une publication dans les 90 jours~?
\item Quels sont les délais médians, en jours, entre la fin d'un trimestre et la publication, puis le 10-Q, et pour un \omterm{def:rs:fundamental-analyst-and-event-features:fiscal}{exercice comptable}, entre la fin de l'exercice et la publication, puis le 10-K~?
\item Quels sont les plus longs délais de dépôt du 10-Q et du 10-K, et comment se comparent-ils aux échéances~?
\item À quelle fréquence le rapport est-il déposé le jour de la publication, et quel est l'écart médian~?
\end{enumerate}

\medskip\noindent\textbf{Partie II --- Les valeurs.}
\begin{enumerate}[resume]
\item Combien de valeurs trimestrielles ont changé après leur premier dépôt, et combien l'ont fait selon un rapport exact~?
\item Quelles ont été les deux divisions d'actions d'Apple, et quel effet ont-elles eu sur son historique par action~?
\item Que s'est-il passé pour le BPA de Microsoft de l'exercice 2017, et pourquoi~?
\item Quel TTM donne l'exercice 2014 d'Apple à partir des trimestres publiés en premier, et quel BPA annuel Apple a-t-il publié~?
\item Comment un entrepôt doit-il conserver les valeurs touchées par une division d'actions~?
\end{enumerate}

\medskip\noindent\textbf{Partie III --- Les dégâts.}
\begin{enumerate}[resume]
\item Quelle est l'avance moyenne d'une clé de fin de période dans \texttt{firm.synthmkt}~?
\item Quels sont les IC de la surprise, indexés sur l'annonce et sur la fin de période, avec leurs erreurs types~?
\item Quelle part des fenêtres contient l'annonce~?
\item Quels sont les IC du ratio livre sur cours sous les deux clés, avec leurs erreurs types~?
\item Pourquoi la différence sur le ratio livre sur cours n'est-elle pas significative~?
\item Énoncer le \emph{résultat nommé}~: la gonflement de l'IC d'une surprise et d'un signal de valeur lorsque la datation à la fin de période remplace la datation à la date connue.
\end{enumerate}

\medskip\noindent\textbf{Partie IV --- Au-delà.}
\begin{enumerate}[resume]
\item Indexer la surprise sur le dépôt plutôt que sur la publication~: que pensez-vous perdre~?
\item Lequel est pire pour un backtest, une clé de fin de période ou un historique de dernière valeur (retraitée), pour une caractéristique SUE~?
\item Comment détecteriez-vous une clé de fin de période dans le backtest d'un autre~?
\item Que doit livrer un fournisseur de fondamentaux pour un usage de recherche~?
\item En une phrase~: quel est le \omterm{def:rs:point-in-time-data-and-the-biases:bitemporal}{temps de connaissance} d'un nombre comptable~?
\end{enumerate}
\end{problem}

\section{Questions d'entretien}

\begin{interviewq}[$\star$ \iqroles{chercheur}]\label{iq:rs:fundamental-analyst-and-event-features:1}
Qu'est-ce que les données fondamentales point-in-time, et pourquoi importent-elles~?
\end{interviewq}

\begin{interviewq}[$\star\star$ \iqroles{chercheur}]\label{iq:rs:fundamental-analyst-and-event-features:2}
Comment calculeriez-vous le résultat des \omterm{def:rs:fundamental-analyst-and-event-features:ttm}{douze derniers mois} (TTM) pour une entreprise qui ne publie que des chiffres cumulés depuis le début de l'exercice~?
\end{interviewq}

\begin{interviewq}[$\star\star$ \iqroles{chercheur, trader}]\label{iq:rs:fundamental-analyst-and-event-features:3}
Qu'est-ce que la dérive post-annonce de résultats, et comment construiriez-vous une caractéristique pour la capter~?
\end{interviewq}

\begin{interviewq}[$\star\star$ \iqroles{chercheur}]\label{iq:rs:fundamental-analyst-and-event-features:4}
Le backtest de votre signal de valeur s'améliore d'un tiers lorsque vous passez de la date de dépôt à la fin de période comme clé. Lequel est juste, et
pourquoi l'amélioration pourrait-elle être faible ou forte~?
\end{interviewq}

\begin{interviewq}[$\star\star$ \iqroles{chercheur}]\label{iq:rs:fundamental-analyst-and-event-features:5}
Quels biais affectent les données de prévisions des analystes~?
\end{interviewq}

\begin{interviewq}[$\star\star\star$ \iqroles{chercheur, développeur}]\label{iq:rs:fundamental-analyst-and-event-features:6}
Concevez le stockage des valeurs comptables par action de sorte que toute requête renvoie ce qui était connu à une date donnée, sur la base d'actions de cette date.
\end{interviewq}
